\section*{ID7512: Quick Numeric Back-of-Envelope (Why Sun Cannot Synchronize Km Bodies at Asteroid: Q∼ 100}
\paragraph{Metadata.}
PiID: 7151249140430 | RTEID: RTE:P29170.7642:T4.0822 | UUID: cc8f456a-d3fa-5da9-997d-aa90650f3a53

\paragraph{Canonical Statement.}
\#\# Quick numeric back-of-envelope (why Sun cannot synchronize km bodies at asteroid/Kuiper distances) Take small body: \textbackslash{}(R\textbackslash{}sim 10\textasciicircum{}3\textbackslash{}) m, \textbackslash{}(m\textbackslash{}sim 10\textasciicircum{}\{12\}\textbackslash{}) kg, typical \textbackslash{}(Q\textbackslash{}sim 100\textbackslash{}), \textbackslash{}(k\_2\textbackslash{}sim 0.1\textbackslash{}). Put these in the \textbackslash{}(\textbackslash{}tau\_\{\textbackslash{}rm sync\}\textbackslash{}) formula with \textbackslash{}(a\textbackslash{}sim 3\textbackslash{}) AU (asteroid belt) or \textbackslash{}(a\textbackslash{}sim 40\textbackslash{}) AU (Kuiper). Because of the \textbackslash{}(a\textasciicircum{}6\textbackslash{}) multiplier, \textbackslash{}(\textbackslash{}tau\_\{\textbackslash{}rm sync\}\textbackslash{}) becomes astronomically large (≫ age of the Solar System) — that’s the core reason.

\paragraph{Canonical Equation.}
\[
Q\sim 100
\]

\paragraph{Dictionary.}
Relation: Q∼ 100

\paragraph{Encyclopedia.}
Quick Numeric Back-of-Envelope (Why Sun Cannot Synchronize Km Bodies at Asteroid: Q∼ 100 is a scaling / order-of-magnitude relation, written Q∼ 100. It expresses a scaling or proportionality, capturing how the quantity grows with its parameters. It is built from Q. Within the Trawin Zero Point registry it serves as one element of the framework's mathematical narrative.
